% Lab part of the module.
%
% Module MC-5 of Section 6.1 of the syllabus: "Keyword spotting features and
% training in Python". Type T+L, 90 minutes: about 45 minutes of theory and
% 45 minutes of lab. The lab needs no board. It runs on a laptop CPU and
% needs the public dataset of Edge Impulse, which a script of the folder
% downloads.
%
% The parts follow the README of the lab folder. The frames show the form of
% the output and the values that check the method. No value here is a
% measurement of a board.

% --------------------------------------------------------------- 1. the goal
\begin{frame}{Lab: goal and plan}
  \small
  \begin{itemize}
    \item \textbf{Goal.} You compute the cepstral coefficients of a clip of
          one second with numpy only, and you train a small classifier on
          them, with no tool of Edge Impulse Studio.
    \item \textbf{Deliverable.} The table of section 10 of the notebook,
          copied into \code{report.md}, and the answers of its three questions.
  \end{itemize}
  \vskip 0.5em
  \footnotesize
  \setlength{\tabcolsep}{4pt}
  \begin{tabular}{@{}L{1.0cm}L{7.2cm}L{1.5cm}@{}}
    \toprule
    Part & Content & Time \\
    \midrule
    A & The clips, and the triangular filters of the Mel scale & 12 min \\
    B & The cepstral coefficients, and the plot of one clip & 13 min \\
    C & The classifier, and its number of parameters & 8 min \\
    D & Three inputs for the same task & 12 min \\
    \bottomrule
  \end{tabular}
  \vskip 0.5em
\end{frame}

% ---------------------------------------------------------- 2. the dataset
\begin{frame}{Lab: the dataset}
  \small
  \begin{itemize}
    \item The script \code{host/get\_keywords.py} downloads the file
          \code{keywords2.zip} of Edge Impulse (139 MB) into the folder
          \code{downloads/}, one time. Then it writes 240 clips of each of
          the four classes into \code{keywords/clips/}.
    \item The four classes are \code{no}, \code{noise}, \code{unknown}, and
          \code{yes}. Each clip has 1 s of audio at 16 kHz.
    \item The notebook reads the amplitudes from text files, so it needs no
          audio package.
  \end{itemize}
  \vskip 0.5em
  \begin{center}
  \begin{tabular}{@{}lrrr@{}}
    \toprule
    Class & Clips & Training & Test \\
    \midrule
    \code{no} & 240 & 192 & 48 \\
    \code{noise} & 240 & 192 & 48 \\
    \code{unknown} & 240 & 192 & 48 \\
    \code{yes} & 240 & 192 & 48 \\
    \bottomrule
  \end{tabular}
  \end{center}
  \small The split is random, so a clip of one speaker can be in both sets.
  The notebook says so in section 6.
\end{frame}

% ------------------------------------------------------------- 3. the files
\begin{frame}{Lab: files}
  \small
  \setlength{\tabcolsep}{4pt}
  \begin{tabular}{@{}L{3.9cm}L{5.3cm}@{}}
    \toprule
    File & Content \\
    \midrule
    \code{mfcc\_training.ipynb} & The notebook. It has four tasks. \\
    \code{host/get\_keywords.py} & Downloads the dataset. \\
    \code{report.md} & The report to hand in. \\
    \code{solutions/} & The notebook with all tasks complete and with its
      output, and an example report. \\
    \code{TEST\_NOTES.md} & The code status and the steps to run the lab. \\
    \bottomrule
  \end{tabular}
  \vskip 0.5em
  \small
  \begin{itemize}
    \item Section 2 of the notebook gives the functions of the pipeline that
          need no decision from you. Two of the four tasks use them.
    \item Section 6 caches the features in the file \code{features.npz}, so
          the second run of the notebook is fast.
  \end{itemize}
\end{frame}

% ---------------------------------------------------------------- 4. part A
\begin{frame}[fragile]{Part A: the clips and the filter bank (12 min)}
  \small
  \begin{enumerate}
    \item Get the dataset, one time, from the folder of the lab:
\begin{lstlisting}[style=shell]
python3 host/get_keywords.py
\end{lstlisting}
    \item Start Jupyter from the folder of the lab, and run sections 0 and 1.
    \item \textbf{Task 1.} Write \code{mel\_filters(n\_filters, n\_bins,
          f\_min, f\_max, sample\_rate)}: the matrix of the 32 triangular
          filters of the Mel scale.
    \item Run the cell after the task.
  \end{enumerate}
  \vskip 0.4em
  \small You see the peak of the first filter at bin 2 and of the last one
  at bin 236. The first filter is 3 bins wide and the last one 40 bins wide.
  That is the Mel scale: narrow at the bottom, wide at the top.
\end{frame}

% ---------------------------------------------------------------- 5. part B
\begin{frame}{Part B: the cepstral coefficients (13 min)}
  \small
  \begin{enumerate}
    \item \textbf{Task 2.} Write \code{mfcc(clip, filters=None)}: the
          pre-emphasis, the framing, the Hamming window, the FFT, the
          filters, the logarithm, and the cosine transform.
    \item Run the cell after the task. Then run section 5.
    \item Look at the three panels. Write in the report what the middle
          panel shows and what the bottom panel shows.
    \item Run section 6. Write the run time of the features in the report.
  \end{enumerate}
  \vskip 0.4em
  \small The notebook prints the mean of the first coefficient, which is
  $-13.071$, and the mean of the last one, which is $0.064$. The first
  coefficient is large because it holds the energy of the frame, and the
  logarithm of a small power is a large negative number.
\end{frame}

% ---------------------------------------------------------------- 6. part C
\begin{frame}{Part C: the classifier (8 min)}
  \small
  \begin{enumerate}
    \item \textbf{Task 3.} Write \code{make\_model(n\_frames, n\_features)}:
          two \code{Conv1D} layers with 16 and with 32 filters of 3 frames,
          each with a max pooling of 2, then a \code{Dense} layer with 4
          outputs.
    \item Run the cell after the task. The check calculates the expected
          number of parameters itself, with the length that each
          convolution and each pooling layer leaves.
    \item Run the next cell to train the model.
  \end{enumerate}
  \vskip 0.5em
  \begin{center}
  \begin{tabular}{@{}l@{}}
    \toprule
    the input has 50 frames of 13 values, and 11 frames after the two \\
    convolutions and the two pooling layers \\
    parameters 3620 (expected 3620), outputs 4 \\
    the cepstral coefficients: 650 values per clip, 3620 parameters, \\
    test accuracy 77.1 percent \\
    \bottomrule
  \end{tabular}
  \end{center}
\end{frame}

% ---------------------------------------------------------------- 7. part D
\begin{frame}{Part D: three inputs for the same task (12 min)}
  \small
  \begin{enumerate}
    \item \textbf{Task 4.} Before you run the next cell, complete
          \code{values\_per\_clip} with the number of values that one clip
          gives for each of the three inputs, and write in
          \code{best\_input} the input that you expect to win.
    \item Run the next cell and section 10, and copy the table into the
          report.
    \item Answer the three questions of the notebook.
  \end{enumerate}
  \vskip 0.4em
  \small The three lines of the notebook: the raw audio has 16000 values
  and 513380 parameters, the log Mel spectrogram 1600 values and 4532
  parameters, and the cepstral coefficients 650 values and 3620 parameters.
\end{frame}

% ------------------------------------------------------------ 8. the check
\begin{frame}{Check criterion and Decision Log}
  \small
  \begin{itemize}
    \item The notebook prints \code{complete} for tasks 1, 2, 3, and 4.
    \item The table of the report has three rows, each with the number of
          values, the number of parameters, and a measured test accuracy.
    \item The three questions of the report have answers with numbers from
          your own run.
    \item The Decision Log names one decision and one number.
  \end{itemize}
  \vskip 0.5em
  \keypoint{Question of the Decision Log: the device must say one word every
    500 ms on a battery. Which of the three inputs do you put in the
    firmware, and which number of your table decides it?}
\end{frame}

% ------------------------------------------------------- 9. common problems
\begin{frame}{Common problems}
  \small
  \setlength{\tabcolsep}{3pt}
  \begin{tabular}{@{}L{4.0cm}L{5.2cm}@{}}
    \toprule
    Problem & Fix \\
    \midrule
    The notebook stops with \code{No clip in the folder} & Run
      \code{host/get\_keywords.py} first. The download is 139 MB. \\
    \code{Task 1} says that the last filter must reach bin 205 & A bin is
      \code{frequency * nfft / sample\_rate}, and \code{nfft} is
      \code{2 * (n\_bins - 1)}. \\
    \code{Task 2} says that the answer needs 50 rows & A clip of 16000
      samples at a stride of 320 gives 50 frames. Return one row for each. \\
    \code{Task 3} says that the model raises an error & A layer is missing,
      or a layer was added with a tensor instead of with its input. \\
    The plot of section 5 is empty & Delete \code{features.npz} and run
      section 6 again. \\
    The notebook is slow & The first run computes the features of 960
      clips. The file \code{features.npz} makes the next run fast. \\
    \bottomrule
  \end{tabular}
  \vskip 0.5em
  \small The README of the lab folder has the full list of the checks of the
  notebook and the run time.
\end{frame}